\chapter{3.30 关于金融支持制造强国建设的指导意见(央行等五部门)} %2017

\begin{center}
中国人民银行~工业和信息化部~银监会~证监会~保监会 \\
关于金融支持制造强国建设的指导意见 \\
银发[2017]58号
\end{center}

为贯彻落实党的十八大和十八届三中、四中、五中、六中全会精神，按照《中国制造2025》、《国务院关于深化制造业与互联网融合发展的指导意见》等要求，进一步建立健全多元化金融服务体系，大力推动金融产品和服务创新，加强和改进对制造强国建设的金融支持和服务，现提出如下意见：

\section{高度重视和持续改进对制造强国建设的金融支持和服务}

（一）坚持问题导向，着力加强对制造业科技创新、转型升级和科技型中小制造企业的金融支持。制造业是实体经济的主体，是科技创新的主战场，是供给侧结构性改革的主攻领域。当前，我国制造业仍存在大而不强，创新能力弱，关键核心技术与高端装备对外依存度高等问题。金融部门应聚焦制造业发展的难点痛点，坚持区别对待、有扶有控原则，不断优化金融支持方向和结构，着力加强对制造业科技创新和技术改造升级的中长期金融支持，积极拓宽技术密集型和中小型制造业企业的多元化融资渠道，促进制造业结构调整、转型升级、提质增效和由大变强。

（二）突出支持重点，改进和完善制造业重点领域和关键任务的金融服务。金融部门要紧紧围绕《中国制造2025》重点任务和“1+X”规划体系，改进和完善金融服务水平。探索为制造业创新中心等公共服务平台提供创新型、多元化融资服务，支持关键共性技术研发和科技成果转化应用。着力加强对核心基础零部件等“四基”企业的融资支持，促进提升工业基础水平。切实加强对企业技术改造中长期贷款支持，合理安排授信期限和还款方式，支持制造业两化融合发展和智能化升级。大力发展绿色金融业务，促进制造业绿色发展。积极运用信贷、租赁、保险等多种金融手段，支持高端装备领域突破发展和扩大应用。加强对工业互联网重点项目的金融支持力度，支持制造业个性化定制和服务型制造。

\section{积极发展和完善支持制造强国建设的多元化金融组织体系}

（三）发挥各类银行机构的差异化优势，形成金融服务协同效应。开发性、政策性金融机构要在重点领域和薄弱环节切实发挥引领作用，在业务范围内以财务可持续为前提，加大对重大项目、重大技术推广和重大装备应用的融资支持。全国性商业银行要积极发挥网络渠道、业务功能协同等优势，为制造业企业提供综合性金融服务，不断改进和提升对中小型制造业企业金融服务的质量效率。地方法人金融机构要注重发挥管理半径短、经营机制灵活等优势，立足当地，服务中小，积极开发针对中小型制造业企业的特色化、专业化金融产品和服务。

（四）完善银行机构组织架构，提升金融服务专业化水平。鼓励有条件的金融机构探索建立先进制造业融资事业部制，加强对信息技术、高端装备、新材料、生物医药等战略重点行业的专业化支持。鼓励符合条件的银行业金融机构在新型工业化产业示范基地等先进制造业聚集地区设立科技金融专营机构，在客户准入、信贷审批、风险偏好、业绩考核等方面实施差异化管理。积极推动小微企业专营机构建设，围绕制造业中量大面广的小微企业、民营企业，提供批量化、规模化、标准化的金融服务。要完善小微企业授信工作尽职免责管理制度，激励基层机构和信贷人员支持中小微制造业企业发展。

（五）规范发展制造业企业集团财务公司。支持符合条件的制造业企业集团设立企业集团财务公司，充分发挥财务公司作为集团“资金归集平台、资金结算平台、资金监控平台、融资营运平台、金融服务平台”的功能，有效提高企业集团内部资金运作效率和精细化管理水平。鼓励具备条件的制造业企业集团财务公司在有效防控风险的前提下，通过开展成员单位产品的买方信贷、消费信贷和融资租赁服务，促进集团产品销售。稳步推进企业集团财务公司开展延伸产业链金融服务试点工作，通过“一头在外”的票据贴现业务和应收账款保理业务，促进降低产业链整体融资成本，更好的支持集团主业发展。

（六）加快制造业领域融资租赁业务发展。积极支持符合条件的金融机构和制造业企业在制造业集聚地区，通过控股、参股等方式发起设立金融租赁公司，支持大型飞机、民用航天、先进轨道交通、海洋工程装备和高技术船舶、智能电网成套设备等高端装备重点领域扩大市场应用和提高国际竞争力。大力发展直接租赁、售后回租等业务，充分发挥融资租赁业务支持企业融资与融物的双重功能，通过“以租代购”、分期偿还等方式，支持制造业企业实施设备更新改造和智能升级。积极发挥融资租赁“以租代售”功能，支持制造业企业扩大销售和出口。

\section{创新发展符合制造业特点的信贷管理体制和金融产品体系}

（七）优化信贷管理体制。鼓励金融机构围绕制造业新型产业链和创新链，积极改进授信评价机制，创新金融产品和服务。合理考量制造业企业技术、人才、市场前景等“软信息”，将相关因素纳入银行客户信用评级体系，挖掘企业潜在价值。鼓励有条件的金融机构在风险可控、商业可持续的前提下，结合企业“三表”、“三单”、“两品”等非财务信息，运用信用贷款、知识产权质押贷款、股权质押贷款、应收账款质押贷款和以品牌为基础的商标专利权质押贷款等方式，积极满足创新型制造业企业和生产性服务业的资金需求。

（八）大力发展产业链金融产品和服务。鼓励金融机构依托制造业产业链核心企业，积极开展仓单质押贷款、应收账款质押贷款、票据贴现、保理、国际国内信用证等各种形式的产业链金融业务，有效满足产业链上下游企业的融资需求。充分发挥人民银行应收账款融资服务平台的公共服务功能，降低银企对接成本。鼓励制造业核心企业、金融机构与人民银行应收账款融资服务平台进行对接，开发全流程、高效率的线上应收账款融资模式。研究推动制造业核心企业在银行间市场注册发行供应链融资票据。

（九）推动投贷联动金融服务模式创新。稳妥有序推进投贷联动业务试点，鼓励和指导试点银行业金融机构以投贷联动方式，为科创型制造业企业提供持续资金支持，促进企业融资结构合理化，有效降低融资成本。建立银行及其投资功能子公司、政府贷款风险补偿基金、融资担保公司、保险公司间的风险分担和补偿机制，有效降低银行信贷风险。鼓励银行业金融机构与外部投资公司、各类基金开展合作，积极整合各自的资金、信息和管理优势，探索多样化的投贷联动业务，促进银企信息交流共享，实现合作共赢。

（十）完善制造业兼并重组的融资服务。推动金融机构对兼并重组企业实行综合授信。鼓励金融机构完善并购贷款业务，在综合考虑并购方的资信状况、经营管理能力、财务稳健性、自筹资本金充足情况，以及并购标的的市场前景、未来盈利、并购协同效应等因素的基础上，合理确定贷款期限和利率，支持企业通过兼并重组实现行业整合。允许符合条件的制造业企业通过发行优先股、可转换债券、并购债券等方式筹集兼并重组资金。鼓励证券公司、资产管理公司、股权投资基金以及产业投资基金等参与企业兼并重组，扩大企业兼并重组资金来源。

（十一）切实择优助强，有效防控风险。银行业金融机构要坚持独立审贷、自主决策、自担风险原则，择优支持有核心竞争力的产业聚集区和企业，注重从源头把控风险。要发挥银行业同业沟通、协调、自律作用，通过联合授信、银团贷款等方式，形成融资协同效应，切实防止多头授信、过度授信，避免一哄而上、重复建设形成新的过剩产能。支持制造业企业按照市场化、法制化原则实施债转股，合理加大股权融资力度，加强企业自身债务杠杆约束，降低企业杠杆率。要稳妥有序退出过剩产能领域，对冶金、建材、石化化工、船舶等行业中有市场、有效益、有技术、经营规范的企业和技改项目，要支持其合理信贷需求。

\section{大力发展多层次资本市场，加强对制造强国建设的资金支持}

（十二）充分发挥股权融资作用。积极支持符合条件的优质、成熟制造业企业在主板市场上市融资，促进重点领域制造业企业做优做强。加快推进高技术制造业企业、先进制造业企业在中小企业板、创业板、全国中小企业股份转让系统和区域性股权交易市场上市或挂牌融资，充实中长期资本实力。在上市融资企业储备库里，对创新能力强、成长性好的制造业企业重点扶持。支持制造业企业在境外上市融资，提升中国制造业企业的国际竞争力。鼓励制造业企业通过资本市场并购重组，实现行业整合和布局调整优化，支持中西部地区承接产业转移。

（十三）促进创业投资持续健康发展。进一步完善扶持创业投资发展的政策体系，促进创业投资发展，有效弥补创新型、成长型制造业企业的融资缺口。鼓励种子基金等各类创业投资基金、天使投资人等创业投资主体加大对种子期、初创期创新型制造业企业的支持力度，通过提供企业管理、商业咨询、财务顾问等多元化服务，支持技术创新完成从科技研发到商业推广的成长历程。鼓励创业投资基金、产业投资基金投向“四基”领域重点项目。发挥先进制造业产业投资基金、国家新兴产业投资引导基金等作用，鼓励建立按市场化方式运作的各类高端装备创新发展基金。

（十四）支持制造业企业发行债券融资。充分发挥公司信用类债券部际协调机制作用，支持符合条件的制造业企业发行公司债、企业债、短期融资券、中期票据、永续票据、定向工具等直接融资工具，拓宽融资渠道，降低融资成本，调整债务结构。设计开发符合先进制造业和战略新兴产业特点的创新债券品种。支持高新技术产业开发区的园区运营机构发行双创专项债务融资工具，用于建设和改造园区基础设施，以及为入园入区制造业企业提供信用增信等服务。支持符合条件的优质制造业企业在注册发行分层分类管理体系下统一注册、自主发行多品种债务融资工具，提升储架发行便利。

（十五）支持制造业领域资产证券化。鼓励金融机构将符合国家产业政策、兼顾收益性和导向性的制造业领域信贷资产作为证券化基础资产，发行信贷资产证券化产品。鼓励制造业企业通过银行间市场发行资产支持票据，以及通过交易所市场开展企业资产证券化，改善企业流动性状况。大力推进高端技术装备、智能制造装备、节能及新能源装备等制造业融资租赁债权资产证券化，拓宽制造业融资租赁机构资金来源，更好服务企业技术升级改造。在依法合规、风险可控的前提下，鼓励符合条件的银行业金融机构稳妥开展不良资产证券化试点，主动化解制造业过剩产能领域信贷风险。

\section{发挥保险市场作用，助推制造业转型升级}

（十六）积极开发促进制造业发展的保险产品。进一步鼓励保险公司发展企业财产保险、科技保险、专利保险、安全生产责任保险等保险业务，为制造业提供多方面的风险保障。鼓励发展制造业贷款保证保险，支持轻资产科创型高技术企业发展壮大。大力发展产品质量责任保险，提高中国制造品牌信任度。深入推进首台（套）重大技术装备保险补偿机制试点工作，推动重大技术装备和关键零部件市场化应用。研究启动重点新材料首批次保险补偿机制试点工作。鼓励地方政府结合本地实际，建立符合本地制造业发展导向的保费补贴和风险补偿机制。

（十七）扩大保险资金对制造业领域投资。积极发挥保险长期资金优势，在符合保险资金运用安全性和收益性的前提下，通过债权、股权、基金、资产支持计划等多种形式，为制造业转型升级提供低成本稳定资金来源。支持保险机构投资制造业企业发行的优先股、并购债券等新型金融工具。鼓励保险机构与银行业金融机构共享信息、优势互补，合作开展制造业领域股债结合、投贷联动等业务。鼓励有条件的保险机构投资设立制造业保险资产管理机构。允许保险资金投资制造业创业投资基金等私募基金。扩大中国保险投资基金对制造业转型升级项目的投入。

\section{拓宽融资渠道，积极支持制造业企业“走出去”}

（十八）拓宽制造业“走出去”的融资渠道。金融机构要根据企业“走出去”需要，进一步优化完善海外机构布局，提高全球化金融服务能力。要积极运用银团贷款、并购贷款、项目融资、出口信贷等多种方式，为制造业企业在境外开展业务活动提供多元化和个性化的金融服务。支持“走出去”企业以境外资产和股权等权益为抵押获得贷款，提高企业融资能力。支持制造业企业开展外汇资金池、跨境双向人民币资金池业务，支持制造业企业在全口径跨境融资宏观审慎管理政策框架下进行跨境融资。支持符合条件的境内制造业企业利用境外市场发行股票、债券和资产证券化产品。

（十九）完善对制造业企业“走出去”的支持政策。不断优化外汇管理，满足制造业企业“走出去”过程中真实、合理的购汇需求。支持制造业企业在对外经济活动中使用人民币计价结算，优化对外人民币贷款项目管理，鼓励企业使用人民币对外贷款和投资。推动设立人民币海外合作基金，为制造业企业“走出去”项目提供成本适当的人民币贷款或投资。鼓励进一步扩大短期出口信用保险规模，加大对中小微企业和新兴市场开拓的保障力度。发挥好中长期出口信用保险的风险保障作用，实现大型成套设备出口融资应保尽保。

\section{加强政策协调和组织保障}

（二十）深入推动产业和金融合作。建立和完善工业和信息化部、人民银行、银监会信息共享和工作联动机制，加强产业政策与金融政策的沟通协调。建立产融信息对接合作平台，促进产业政策信息、企业生产经营信息、金融产品信息交流共享。探索对制造业部分重点行业建立企业“白名单”制度，为金融机构落实差别化的信贷政策提供参考依据。开展产融合作城市试点，促进试点城市聚合产业资源、金融资源、政策资源，支持制造强国建设，鼓励和引导产融合作试点城市探索产业与金融良性互动、实现互利共赢。

（二十一）加大货币信贷政策的支持力度。保持货币政策稳健中性，综合运用多种流动性管理工具，加强预调微调，保持流动性水平合理适度，引导货币信贷平稳增长，营造良好的货币金融环境，为制造业企业加强技术创新、实现转型升级创造有利的外部条件。充分发挥再贷款、再贴现等货币政策工具的正向激励作用，对支持制造业转型升级工作成绩突出的金融机构要优先予以支持。进一步引导银行业金融机构完善存贷款利率定价机制，增强自主理性定价能力，合理确定对制造业企业的贷款利率水平。

（二十二）优化政策配套和协调配合。鼓励地方政府通过设立产业引导基金、股权投资引导基金，按照政府引导、市场化运作、专业化管理的原则，加大对先进制造业的投资力度，撬动各类社会资本支持制造业转型升级。积极探索多样化的信贷风险分担机制，通过设立风险补偿基金、政府性担保基金、应急转贷基金等方式，支持金融机构加大制造业领域的信贷投入。进一步完善银行与融资担保机构的合作机制，建立合理的企业贷款风险分担和利率协商机制。进一步加强社会信用体系建设，对恶意逃废债行为要给予严厉打击，维护金融机构合法权益。

（二十三）加强监督引导和统计监测。建立金融支持制造强国建设工作的部际协调机制，加强信息沟通和监督引导。人民银行分支机构、银监会、证监会、保监会派出机构要根据当地产业发展实际情况，适时研究制定配套金融政策措施，切实加强对金融机构的指导。要依托行业主管部门支持，研究建立支持制造强国建设的金融统计制度，进一步加强对高技术制造业、先进制造业融资情况的统计监测分析。

请人民银行上海总部，各分行、营业管理部、省会（首府）城市中心支行、副省级城市中心支行会同所在省（区、市）工业和信息化主管部门、银监局、证监局、保监局等部门，将本意见联合转发至辖区内相关机构，并协调做好本意见的贯彻实施工作。
